% Standalone mathematical input: uses only amsmath, amssymb, amsthm
% and the containing article's theorem/proof environments.
% Citation key supplied by the containing article: ConnesHeat.

\section{Gaussian heat in the original source, full jets and attained metrics}

We use the Gaussian in Alain Connes, \emph{Heat expansion and zeta},
Sections 2--3, original source labels \texttt{mellin},
\texttt{fourierm}, \texttt{ft} and \texttt{plicit1}
\cite{ConnesHeat}.  His multiplicative test function and its Mellin
version are, with every factor retained,
\begin{equation}\tag{GM1}
 K_t(v)=\frac{e^{-v^2/(4t)}}{\sqrt{4\pi t}},\qquad
 F_t(x)=K_t(\log x),\qquad f_t(x)=x^{-1/2}F_t(x),\quad t>0.
\end{equation}
Connes uses $\widehat F(s)=\int_0^\infty F(x)x^{-is}\,dx/x$.
Consequently $\widehat F_t(s)=e^{-ts^2}$ and
$\mathcal M f_t(w)=e^{t(w-1/2)^2}$.  The source Hilbert space below
has the original measure $dx$, not $dx/x$.

\subsection{The source operator and the complete physical jet unit}

Put $\mathcal H=L^2((0,\infty),dx)$, $D=-x\partial_x$, and
\[
 (UF)(v)=e^{v/2}F(e^v),\qquad
 \mathcal M F(w)=\int_0^\infty F(x)x^{w-1}\,dx.
\]
The smooth source space used here is
\[
 \mathscr B=\{F\in C^\infty(0,\infty):
       \sup_{x>0}x^a|D^bF(x)|<\infty
       \text{ for every }a\in\mathbb R, b\geq0\}.
\]
Thus $UF$ and each of its derivatives decrease faster than every
exponential at both ends of the real line.  Keep
$g(w)=2\xi(w)=w(w-1)\pi^{-w/2}\Gamma(w/2)\zeta(w)$ and a full finite
divisor $h(w)=\prod_{\rho\in Z}(w-\rho)^{m_\rho}$, where $m_\rho$
is the actual multiplicity of $g$ at $\rho$.  Set $d=\deg h$.
The source $a_h\in\mathscr B$ is characterized by
$\mathcal M a_h=g/h$.  This characterization preserves the arithmetic
unit at every root.  Its construction is the following exact division
of the original theta source:
\[
 f_0(x)=2\sum_{n\geq1}(4\pi^2 n^4x^4-6\pi n^2x^2)e^{-\pi n^2x^2},
 \qquad \mathcal M f_0=g,
 \qquad
 S_\rho F(x)=x^{-\rho}\int_x^\infty F(z)z^{\rho-1}\,dz.
\]
When $\mathcal M F(\rho)=0$, the last integral also equals
$-x^{-\rho}\int_0^xF(z)z^{\rho-1}\,dz$.  These two expressions show
rapid decrease at the respective endpoints, including every derivative.
Differentiating gives $(D-\rho)S_\rho F=F$; integration by parts gives
$\mathcal M(S_\rho F)=\mathcal M F/(w-\rho)$ with its removable value.
Applying this operation to the complete multiset of roots produces
$a_h$.  Mellin injectivity proves its uniqueness.  For completeness,
the displayed identity for $f_0$ follows by integrating its absolutely
convergent theta sum for $\Re w>1$, using
$D(D-1)e^{-\pi x^2}=(4\pi^2x^4-6\pi x^2)e^{-\pi x^2}$,
and then continuing.  Poisson summation for
$\sum_{n\in\mathbb Z}e^{-\pi n^2x^2}$ shows that $f_0$ and every
$D$ derivative decrease faster than every power at zero as well as
infinity: the two nondecreasing terms $x^{-1}$ and $1$ are annihilated
by $D(D-1)$.  This also verifies $f_0\in\mathscr B$.

\begin{theorem}[GM-A: exact Gaussian source and full jet morphism]
The contraction semigroup $T(t)$ on $\mathcal H$ defined by
$UT(t)U^{-1}=K_t*$ has generator $(D-1/2)^2$ and satisfies
\begin{equation}\tag{GM2}
 \begin{split}
 T(t)F(x)&=\int_0^\infty
       x^{-1/2}z^{-1/2}K_t(\log(x/z))F(z)\,dz,\\
 \mathcal M(T(t)F)(w)&=e^{t(w-1/2)^2}\mathcal M F(w)
       \quad(F\in\mathscr B, w\in\mathbb C),\\
 \|T(t)P(D)a_h\|^2&=\int_{\mathbb R}
       |P(1/2+iy)|^2 e^{-2ty^2} w_h(y)\,dy,\qquad
 w_h(y)=\frac{|(g/h)(1/2+iy)|^2}{2\pi}.
 \end{split}
\end{equation}
The map $T(t)$ preserves $\mathscr B$ and is injective.  Let
$E_h=\prod_{\rho\in Z}\mathbb C[\tau_\rho]/(\tau_\rho^{m_\rho})$
in the literal coordinates $\tau_\rho=w-\rho$, and let $j_h$ take
all divided Taylor coefficients.  Define $B_h(t)$ to be multiplication
in $E_h$ by $j_h(e^{t(w-1/2)^2})$.  Then
\begin{equation}\tag{GM3}
 \begin{gathered}
 j_h\mathcal M(T(t)F)=B_h(t)j_h\mathcal M F,\qquad
 B_h(t)=e^{tL_h},\quad L_h=j_h((w-1/2)^2)\cdot,\\
 (B_h(t))_{(\rho,r),(\rho,s)}
 =e^{t(\rho-1/2)^2}
 \sum_{j=0}^{\lfloor(r-s)/2\rfloor}
 \frac{(2t(\rho-1/2))^{r-s-2j}t^j}{(r-s-2j)!\,j!}
 \quad(r\geq s).
 \end{gathered}
\end{equation}
Entries with $r<s$ or different roots are zero.  In particular this is
an invertible full divided-jet map, including all nilpotent terms.
\end{theorem}
\begin{proof}
The substitution $x=e^v$ proves that $U$ is unitary and that
$U(D-1/2)U^{-1}=-\partial_v$.  Fourier transformation with the
positive exponential $e^{iyv}$ sends this operator to multiplication
by $iy$, and sends $K_t*$ to multiplication by $e^{-ty^2}$.
This proves the generator, contraction, semigroup and injectivity
statements, with the closure of the generator understood on its
Fourier domain.  Substitution $z=e^u$ in the convolution proves its
kernel against the original measure $dz$.

If $|\partial_v^b UF(v)|\leq C_{A,b}e^{-A|v|}$, convolution is bounded
by $C_{A,b}e^{-A|v|}\int K_t(u)e^{A|u|}\,du$; the integral is finite.
This applies to every $A,b$, proves preservation of $\mathscr B$ and
justifies differentiation and all complex Mellin integrals.  Fubini
then gives their multiplier $\int K_t(v)e^{(w-1/2)v}\,dv
=e^{t(w-1/2)^2}$.  The Gaussian integral is first obtained by
completing the square for real $w-1/2$ and then by entire continuation.
Plancherel in the same convention gives
$\|F\|^2=(2\pi)^{-1}\int|\mathcal M F(1/2+iy)|^2dy$, proving the
last identity of (GM2).  Finally expand
$e^{t(\rho-1/2+\tau)^2}
=e^{t(\rho-1/2)^2}e^{2t(\rho-1/2)\tau}e^{t\tau^2}$
modulo $\tau^{m_\rho}$.  This proves every entry in (GM3), and the
inverse is the same finite algebra exponential at $-t$.
\end{proof}

Fix $N\geq d-1$.  For $u\in E_h$ let $r_h(u)$ be the unique
polynomial of degree less than $d$ such that
$j_h r_h(u)=j_h(g/h)^{-1}u$.  The inverse exists because the
multiplicities in $h$ are full.  Hermite interpolation here is simply
the Chinese remainder map $\mathbb C[w]/(h)\simeq E_h$; its kernel
consists exactly of polynomials divisible by every
$(w-\rho)^{m_\rho}$, hence by $h$, and equal finite dimensions give
the asserted isomorphism.  Put $\mathcal P_n=\mathbb C[w]_{\leq n}$,
with $\mathcal P_{-1}=0$, and keep the coefficient map
\begin{equation}\tag{GM4}
 \begin{gathered}
 L_N(u,a)=r_h(u)+ha,\quad
 \mathcal R_N(u,a)=L_N(u,a)(D)a_h,\qquad
 (u,a)\in E_h\oplus\mathcal P_{N-d},\\
 j_h\mathcal M\mathcal R_N(u,a)=u,\qquad
 j_h\mathcal M T(t)\mathcal R_N(u,a)=B_h(t)u.
 \end{gathered}
\end{equation}
Euclidean division by $h$ proves that $L_N$ is an isomorphism.  All
matrices below use these jet coordinates and the ordinary power
coefficients of $a$.  No relation coefficient is discarded.

\subsection{All moment derivatives and the moving attained minimum}

Let $b_N(y)$ be the row vector specified by
$b_N(y)(u,a)=L_N(u,a)(1/2+iy)$.  For $r\geq0$ set
\begin{equation}\tag{GM5}
 \begin{gathered}
 M_r(t)=\int_{\mathbb R}y^{2r}e^{-2ty^2}
                       b_N(y)^*b_N(y)w_h(y)\,dy,\qquad
 G(t)=M_0(t),\\
 G(t)=(T(t)\mathcal R_N)^*T(t)\mathcal R_N
       =\begin{pmatrix}A(t)&C(t)\\ C(t)^*&S(t)\end{pmatrix}.
 \end{gathered}
\end{equation}
All these moments are finite: the Mellin transform of $a_h\in\mathscr B$
decreases faster than every inverse power on the line in (GM5), by
integration by parts in $v$.  Moreover $G(t)>0$, since the entire
nonzero function $g/h$ cannot vanish on an interval and a polynomial
vanishing almost everywhere on that line is zero.

\begin{theorem}[GM-B: moment derivatives with a signed finite remainder]
For every integer $r\geq0$ and $t\geq0$, with right derivatives at zero,
\begin{equation}\tag{GM6}
 G^{(r)}(t)=(-2)^rM_r(t).
\end{equation}
For every $m\geq0$,
\begin{equation}\tag{GM7}
 0\preceq(-1)^{m+1}
 \left[G(t)-\sum_{j=0}^m\frac{(-2t)^j}{j!}M_j(0)\right]
 \preceq\frac{(2t)^{m+1}}{(m+1)!}M_{m+1}(0).
\end{equation}
The exact finite constants
\begin{equation}\tag{GM8}
 \gamma_N=\|G(0)^{-1/2}M_1(0)G(0)^{-1/2}\|,
 \qquad
 \beta_N=\|G(0)^{-1/2}M_2(0)G(0)^{-1/2}\|
\end{equation}
satisfy, for all $t\geq0$,
\begin{equation}\tag{GM9}
 e^{-2t\gamma_N}G(0)\preceq G(t)\preceq G(0).
\end{equation}
\end{theorem}
\begin{proof}
Each differentiated entry is bounded by an integrable polynomial
moment, proving (GM6) by dominated convergence.  For $x\geq0$ the
integral Taylor remainder of $e^{-x}$ after degree $m$ has sign
$(-1)^{m+1}$ and absolute value at most $x^{m+1}/(m+1)!$.
Multiply this scalar inequality, with $x=2ty^2$, by
$|b_N(y)c|^2w_h(y)$ and integrate for each coefficient vector $c$.
This proves (GM7) as an inequality of Hermitian forms.
For $c\ne0$, divide this same positive density at $t=0$ by $c^*G(0)c$
to obtain a probability measure.  Its $y^2$ expectation is at most
$\gamma_N$.  Convexity of $v\mapsto e^{-2tv}$ gives
$c^*G(t)c/(c^*G(0)c)\geq e^{-2t\gamma_N}$; the upper bound follows
from $e^{-2ty^2}\leq1$.  These are precisely (GM9).
\end{proof}

To include all attained quotients, let $J$ be any fixed surjection
from this coefficient space onto a finite coordinate space.  Let $K$
be any fixed injective matrix whose image is $\ker J$; empty matrices
are allowed when that kernel is zero.  Define
\begin{equation}\tag{GM10}
 \begin{gathered}
 H(t)=(JG(t)^{-1}J^*)^{-1},\qquad
 L(t)=G(t)^{-1}J^*H(t),\qquad S_K(t)=K^*G(t)K,\\
 u^*H(t)u=\min_{Jc=u}c^*G(t)c,\qquad JL(t)=I,
 \qquad K^*G(t)L(t)=0.
 \end{gathered}
\end{equation}
Indeed the stated lift is feasible, and every other lift is $L(t)u+Ka$.
Its norm squared is $u^*H(t)u+a^*S_K(t)a$.  This proves the attained
minimum and its uniqueness, including every relation variable.

\begin{theorem}[GM-C: derivative and Hessian of the full attained form]
The lift and minimum in (GM10) obey
\begin{equation}\tag{GM11}
 \begin{split}
 L'&=-K S_K^{-1}K^*G'L,\\
 H'&=L^*G'L=-2L^*M_1L,\\
 H''&=L^*G''L-2L^*G'K S_K^{-1}K^*G'L\\
     &=4L^*M_2L-8L^*M_1K S_K^{-1}K^*M_1L.
 \end{split}
\end{equation}
All matrices on a right-hand side have the same argument $t$.
For the original jet projection $J(u,a)=u$, these formulas are the
full Schur identities
\begin{equation}\tag{GM12}
 \begin{gathered}
 Z=-S^{-1}C^*,\quad L=\binom{I}{Z},\quad H=A-CS^{-1}C^*,\\
 H'=A'-C'S^{-1}C^*-CS^{-1}(C')^*
                       +CS^{-1}S'S^{-1}C^*,\\
 H''=L^*G''L
     -2\bigl(C'+Z^*S'\bigr)S^{-1}\bigl((C')^*+S'Z\bigr).
 \end{gathered}
\end{equation}
When the relation space is zero the inverse terms are absent.
In particular neither cross block nor the motion of its minimizing
coefficient is suppressed.
\end{theorem}
\begin{proof}
Differentiate $JL=I$ to see that every column of $L'$ lies in
$\operatorname{im}K$.  Differentiating $K^*GL=0$ then gives the first
formula in (GM11).  Differentiating $H=L^*GL$ makes the terms
$L'^*GL$ and $L^*GL'$ vanish, by the same orthogonality.  A second
derivative gives $L'^*G'L+L^*G''L+L^*G'L'$.  Substitution of $L'$
and Hermitian symmetry of $G'$ gives exactly the two correction terms
combined in (GM11).  Equation (GM6) gives the last expression.
The block choice $K=\binom0I$ proves (GM12) by multiplication;
alternatively differentiating $Z=-S^{-1}C^*$ gives
$Z'=-S^{-1}((C')^*+S'Z)$ and the same Hessian.
\end{proof}

\begin{theorem}[GM-D: attained finite error and complete physical jets]
For every quotient (GM10),
\begin{equation}\tag{GM13}
 \begin{gathered}
 e^{-2t\gamma_N}H(0)\preceq H(t)\preceq H(0),\\
 -4L(t)^*M_2(t)L(t)\preceq H''(t)
                     \preceq4L(t)^*M_2(t)L(t),\\
 -R_N(t)H(0)\preceq H(t)-H(0)-tH'(0)\preceq R_N(t)H(0),\\
 R_N(t)=4\beta_N\int_0^t(t-s)e^{2s\gamma_N}\,ds
 \leq2\beta_Nt^2e^{2t\gamma_N}.
 \end{gathered}
\end{equation}
For $\gamma_N>0$ the exact displayed remainder constant is
$R_N(t)=\beta_N(e^{2t\gamma_N}-1-2t\gamma_N)/\gamma_N^2$.
For the original projection onto $E_h$, the metric at a fixed
\emph{physical} smoothed jet $v$ is
\begin{equation}\tag{GM14}
 \begin{split}
 v^*\widehat H(t)v
 &=\min_{j_h\mathcal M T(t)\mathcal R_Nc=v}
                  \|T(t)\mathcal R_Nc\|^2,\\
 \widehat H(t)&=B_h(t)^{-*}H(t)B_h(t)^{-1},\\
 \widehat H'(t)&=B_h(t)^{-*}
       [H'-L_h^*H-HL_h]B_h(t)^{-1},\\
 \widehat H''(t)&=B_h(t)^{-*}
       [H''-2L_h^*H'-2H'L_h+(L_h^*)^2H
                    +2L_h^*HL_h+HL_h^2]B_h(t)^{-1}.
 \end{split}
\end{equation}
In particular, with $\ell_N=\|H(0)^{1/2}L_hH(0)^{-1/2}\|$,
\begin{equation}\tag{GM15}
 \begin{gathered}
 e^{-2t(\gamma_N+\ell_N)}H(0)
       \preceq\widehat H(t)\preceq e^{2t\ell_N}H(0),\\
 -\widehat R_N(t)H(0)
 \preceq\widehat H(t)-H(0)-t\widehat H'(0)
 \preceq\widehat R_N(t)H(0),\\
 \widehat R_N(t)=2t^2
       (\beta_N+2\ell_N\gamma_N+\ell_N^2)
                     e^{2t(\gamma_N+\ell_N)}.
 \end{gathered}
\end{equation}
\end{theorem}
\begin{proof}
Take the minimum of (GM9) on the identical affine fibre $Jc=u$.
For the Hessian bound, realize the coefficient map as the functions
$c\mapsto e^{-ty^2}b_N(y)c$ in $L^2(w_h(y)dy)$.  Its restriction to
$K$ has Gram $S_K$.  The orthogonal projection onto that finite
relation image is therefore this restricted map times $S_K^{-1}$
times its adjoint.  Apply the contraction bound for this projection
to the function $y^2e^{-ty^2}b_N(y)L(t)u$.  It gives, for every $u$,
\[
 u^*L^*M_1K S_K^{-1}K^*M_1Lu\leq u^*L^*M_2Lu.
\]
Combining it with (GM11) gives the Hessian inequalities in (GM13).
Moreover
\[
 M_2(t)\preceq M_2(0)\preceq\beta_NG(0)
       \preceq\beta_Ne^{2t\gamma_N}G(t),
 \qquad H(t)\preceq H(0).
\]
Thus $L(t)^*M_2(t)L(t)\preceq
\beta_Ne^{2t\gamma_N}H(0)$.  Integrating the Hessian bound twice
proves the remainder, including its integral constant.

Equation (GM4) says exactly that the feasible coefficient label for
physical jet $v$ is $u=B_h(t)^{-1}v$.  Substitution into (GM10) proves
the first two lines of (GM14).  Since $B_h'=L_hB_h=B_hL_h$,
ordinary differentiation proves the remaining lines.

To check (GM15), perform all estimates in the fixed norm
$\|u\|_{H(0)}^2=u^*H(0)u$.  The power series and submultiplicativity
give $\|e^{\pm tL_h}\|_{H(0)}\leq e^{t\ell_N}$.  In particular
$\|e^{-tL_h}u\|_{H(0)}\geq e^{-t\ell_N}\|u\|_{H(0)}$ by applying
the bound for the inverse.  This and (GM13) prove the first line.
The same comparison with $M_1$ in place of $M_2$ gives
$\|H(0)^{-1/2}H'(t)H(0)^{-1/2}\|
\leq2\gamma_Ne^{2t\gamma_N}$.  The Hessian estimate already gives
the analogous bound $4\beta_Ne^{2t\gamma_N}$ for $H''$; also
$\|H(0)^{-1/2}H(t)H(0)^{-1/2}\|\leq1$.
Consequently the bracket in the last line of (GM14) has norm at most
$4(\beta_N+2\ell_N\gamma_N+\ell_N^2)e^{2t\gamma_N}$ in these fixed
coordinates.  Congruence by $B_h^{-1}$ adds at most $e^{2t\ell_N}$.
Integrating twice and bounding the exponential by its endpoint
proves the stated $\widehat R_N$.
\end{proof}

These estimates also cover any specified finite sequence of fixed
restrictions and attained quotients.  Here is the exact reason.
A fixed restriction replaces the source coefficient space by a fixed
linear subspace.  A quotient minimizes over a fixed affine fibre in
that subspace.  Restricting its target to a further fixed subspace
restricts the preceding source to its inverse image; taking a further
quotient composes the fixed surjections.  Thus, by induction, each
resulting metric is a single minimum of the original form on a fixed
subspace with a fixed surjection.  Apply the proofs above to that
subspace.  The inequalities $M_j(0)\preceq c_jG(0)$ restrict with the
same constants, so (GM13) still holds with $\gamma_N,\beta_N$ from
the full original source.  Metric-dependent changes of frame are
not used in this statement.

For a rank-$r$ result $Q(t)$ in such fixed coordinates, (GM13) implies
\begin{equation}\tag{GM16}
 -2rt\gamma_N\leq\log\frac{\det Q(t)}{\det Q(0)}\leq0.
\end{equation}
Indeed conjugation by $Q(0)^{-1/2}$ puts all its eigenvalues in
$[e^{-2t\gamma_N},1]$, and their logarithms add.  Rank zero has
determinant $1$.  Therefore a finite specified return
$\mathcal L(t)=\sum_\nu a_\nu\log\det Q_\nu(t)$ obeys
$|\mathcal L(t)-\mathcal L(0)|
\leq2t\sum_\nu|a_\nu|r_\nu\gamma_{N_\nu}$.
All endpoint coefficients and ranks remain in this error.  For
physical jets the exact additional determinant term is
\begin{equation}\tag{GM17}
 \log\frac{\det\widehat H(t)}{\det H(0)}
 =\log\frac{\det H(t)}{\det H(0)}
       -2t\Re\sum_{\rho\in Z}m_\rho(\rho-1/2)^2.
\end{equation}
This follows because $L_h$ is triangular with diagonal entry
$(\rho-1/2)^2$ repeated $m_\rho$ times.  It is an exact unit
contribution, not an omitted scalar normalization.

\subsection{Tensor heat, sum heat and the original complete quartet}

Retain $k$ separate physical variables, $D_j=-x_j\partial_{x_j}$,
and $\mathcal H^{\otimes k}=L^2((0,\infty)^k,dx_1\cdots dx_k)$.
Write $A_j=D_j-1/2$ and $S=\sum_jw_j$, with
$w_j=1/2+iy_j$.  Define the three commuting contraction semigroups
by their Fourier multipliers:
\begin{equation}\tag{GM18}
 \begin{array}{c|c|c}
 &\text{generator}&\text{multiplier at }(y_1,\ldots,y_k)\\ \hline
 T_{\rm fac}(t)&\sum_jA_j^2&e^{-t\sum_jy_j^2}\\
 T_{\rm sum}(t)&(\sum_jA_j)^2&e^{-t(\sum_jy_j)^2}\\
 T_{\perp}(t)&\sum_jA_j^2-k^{-1}(\sum_jA_j)^2
                &e^{-t\sum_j(y_j-k^{-1}\sum_i y_i)^2}.
 \end{array}
\end{equation}
In particular $T_{\rm fac}(t)=T(t)^{\otimes k}$.  The original
additive source is
\[
 \mathcal A_{k,N}P
 =P(D_1+\cdots+D_k)(a_h\otimes\cdots\otimes a_h),
 \qquad P\in\mathbb C[S]_{\leq N}.
\]
Its multiple Mellin transform is exactly
$P(S)\prod_j(g/h)(w_j)$; this follows by expanding each monomial
with its original multinomial coefficients and using integration by
parts in every variable.

\begin{theorem}[GM-E: exact tensor-to-sum heat map and moment receiver]
There is the bounded operator identity
\begin{equation}\tag{GM19}
 T_{\rm fac}(t)=T_{\rm sum}(t/k)T_{\perp}(t).
\end{equation}
The complete Mellin multipliers of the smoothed additive sources are
respectively
\begin{equation}\tag{GM20}
 e^{t\sum_j(w_j-1/2)^2}P(S)\prod_j(g/h)(w_j),\qquad
 e^{t(S-k/2)^2}P(S)\prod_j(g/h)(w_j).
\end{equation}
Their original power-coefficient Gram forms are
\begin{equation}\tag{GM21}
 \begin{split}
 P^*\mathcal G_{\rm fac}(t)P
 &=\int_{\mathbb R}|P(k/2+iy)|^2
               (e^{-2t(\cdot)^2}w_h)^{*k}(y)\,dy,\\
 P^*\mathcal G_{\rm sum}(t)P
 &=\int_{\mathbb R}|P(k/2+iy)|^2
                 e^{-2ty^2}w_h^{*k}(y)\,dy,\\
 \mathcal G_{\rm fac}(t)&\preceq\mathcal G_{\rm sum}(t/k).
 \end{split}
\end{equation}
All of (GM6)--(GM13) apply to each of these two full Grams and any
of its fixed attained quotients, with their respective original
moments
\begin{equation}\tag{GM22}
 M_r^\diamond(t)=\int_{\mathbb R^k}
   \eta_\diamond(\mathbf y)^r e^{-2t\eta_\diamond(\mathbf y)}
   b_{k,N}(\mathbf y)^*b_{k,N}(\mathbf y)
                          \prod_jw_h(y_j)\,d\mathbf y,
 \quad
 \eta_{\rm fac}=\sum_jy_j^2,\quad
 \eta_{\rm sum}=(\sum_jy_j)^2,
\end{equation}
where $b_{k,N}(\mathbf y)$ is the row evaluating the original
$S$-power basis at $S=k/2+i\sum_jy_j$.
\end{theorem}
\begin{proof}
The identity
$\sum_jy_j^2=k^{-1}(\sum_jy_j)^2+
\sum_j(y_j-k^{-1}\sum_i y_i)^2$ proves (GM19) pointwise on the
unitary multiple Fourier transform.  Each factor is bounded, so this
proves an equality on the whole Hilbert space, with no inverse heat
operator required.  The entire Mellin version gives (GM20), first on
tensor products of $\mathscr B$ and their finite derivative sums,
which include the displayed sources.  The same Gaussian convolution
argument as in GM-A justifies this at all complex coordinates.
Multiple Plancherel gives the $k$-fold integrals of
$|P(k/2+i\sum y_j)|^2$ against the two exponential densities.
Integrate with $y=\sum y_j$ and
$y_k=y-\sum_{j<k}y_j$.  This change of variables has determinant
one, and Tonelli gives the two convolutions in (GM21).  The remaining
factor $e^{-2t\sum_j(y_j-y/k)^2}\leq1$ proves their comparison.
Every polynomial moment in (GM22) is finite, since it is bounded by
a finite sum of products of the one-factor moments.  Positivity
follows from the same polynomial argument as for (GM5).
The proofs of GM-B--D use only a nonnegative measurable multiplier
$\eta$, its moments, and projection onto the actual relation space.
Replacing $y^2$ by the explicitly displayed $\eta_\diamond$ proves
each formula, without changing its factors of two or its minima.
\end{proof}

For example the full tensor coefficient space
$\mathcal P_N^{\otimes k}$ gives valid bounds
$\gamma_{\rm fac}\leq k\gamma_N$ and
$\gamma_{\rm sum}\leq k^2\gamma_N$ after its exact restriction to
the additive source.  Here $\gamma_N$ can be computed in either
original coefficient frame, since its definition is the supremum
of the same Rayleigh quotient.  Indeed each insertion of the
one-factor first moment is at most $\gamma_N$ times the full tensor
Gram, and $(\sum_jy_j)^2\leq k\sum_jy_j^2$.  Likewise
$\beta_{\rm fac}\leq k^2\beta_N$ and
$\beta_{\rm sum}\leq k^4\beta_N$, using
$(\sum_jy_j^2)^2\leq k\sum_jy_j^4$.  These inequalities explicitly
retain both $k$ and the degree-dependent original constants.

Now specialize to the original quartet, writing its positive
displacement as $\delta_\rho$ throughout:
\begin{equation}\tag{GM23}
 \begin{gathered}
 \rho=\tfrac12+\delta_\rho+i\gamma,\quad
 0<\delta_\rho<\tfrac12,\quad\gamma>2,\quad m_0\geq1,\qquad
 k\equiv1\pmod4,\quad k\geq5,\\
 h(w)=\prod_{\epsilon,\eta\in\{1,-1\}}
       [w-\tfrac12-\epsilon\delta_\rho-i\eta\gamma]^{m_0},
 \quad e=1+k(m_0-1),\quad q=e(k+1)^2,\\
 \kappa_{a,b}=k/2+(2a-k)\delta_\rho+i(2b-k)\gamma,\qquad
 \chi(S)=\prod_{a,b=0}^k(S-\kappa_{a,b})^e.
 \end{gathered}
\end{equation}
This calculation studies the stipulated off-axis quartet with its
given parameters; it does not assert that zeta has such a quartet.
Let $\mathcal T_h=E_h^{\otimes k}$ and let $s=\sum_jw_j$ in this
finite algebra.  Multiplication by the full arithmetic unit
$\prod_j j_h(g/h)(w_j)$ is denoted $\mathsf U_h$.

\begin{theorem}[GM-F: complete sum embedding and exact descent exceptions]
There is an injective algebra map
\begin{equation}\tag{GM24}
 \iota:\ E_\chi=\mathbb C[S]/(\chi)\longrightarrow\mathcal T_h,
                 \qquad[P]\longmapsto P(s).
\end{equation}
The physical full jet maps of the three additive sources (unheated,
factor heat, sum heat) are exactly
\begin{equation}\tag{GM25}
 \begin{gathered}
 \mathsf U_h\iota([P]),\qquad
 \mathsf U_h Q_t\iota([P]),\qquad
 \mathsf U_h\iota(E_t[P]),\\
 Q_t=\exp\!\left(t\sum_j(w_j-1/2)^2\right),\quad
 E_t=\exp\bigl(t(S-k/2)^2\bigr)\in E_\chi.
 \end{gathered}
\end{equation}
Finite algebra exponentials here retain all nilpotent coefficients.
For $t>0$, the factor unit $Q_t$ belongs to $\iota(E_\chi)$ exactly
in the following case:
\begin{equation}\tag{GM26}
 m_0=1\quad\hbox{and}\quad4t\delta_\rho\gamma\in\pi\mathbb Z.
\end{equation}
For $m_0\geq2$ it never belongs to that subalgebra at positive time.
For every $m_0,t$ the multiplication map
\begin{equation}\tag{GM27}
 \mathsf U_h\iota(E_\chi)\longrightarrow
                  \mathsf U_h Q_t\iota(E_\chi),\qquad v\longmapsto Q_tv
\end{equation}
is an explicitly invertible linear map, with inverse multiplication
by $Q_{-t}$.  Thus failure of descent does not remove the exact
correspondence between the original physical jet images.
\end{theorem}
\begin{proof}
Decompose $\mathcal T_h$ into its primary factors, one for each ordered
tuple of quartet roots.  At such a tuple write
$w_j=1/2+\epsilon_j\delta_\rho+i\eta_j\gamma+\tau_j$,
with $\tau_j^{m_0}=0$.  The factor is
$\mathbb C[\tau_1,\ldots,\tau_k]/(\tau_1^{m_0},\ldots,\tau_k^{m_0})$.
If precisely $a$ of the $\epsilon_j$ and $b$ of the $\eta_j$ equal
$1$, then $s=\kappa_{a,b}+\sum_j\tau_j$.  The latter nilpotent sum
has index exactly $e$: its $e$-th power is zero by degree, while
\[
 (\textstyle\sum_j\tau_j)^{e-1}
 =\frac{(k(m_0-1))!}{((m_0-1)!)^k}\prod_j\tau_j^{m_0-1}\ne0.
\]
Every pair $(a,b)$ occurs, by choosing independently the sets of
positive $\epsilon_j$ and positive $\eta_j$.  The $\kappa_{a,b}$
are distinct, since their real and imaginary differences vanish only
when both indices agree.  The minimal polynomial of $s$ on the
whole product is therefore exactly $\chi$, proving (GM24) including
every primary multiplicity.  Taking divided Taylor coefficients
of (GM20) proves (GM25); the factor $\mathsf U_h$ is a unit because
every factor $j_h(g/h)$ is a unit.

First take $m_0=1$.  A function on the ordered root tuples is in
$\iota(E_\chi)$ precisely when it is constant on tuples with the
same pair $(a,b)$: necessity follows from evaluation at $s$, and
sufficiency follows by polynomial interpolation at the distinct
$\kappa_{a,b}$.  If $c$ is the number of indices with
$\epsilon_j=\eta_j=1$, direct counting gives
\[
 \sum_j(\epsilon_j\delta_\rho+i\eta_j\gamma)^2
 =k(\delta_\rho^2-\gamma^2)
        +2i\delta_\rho\gamma(k-2a-2b+4c).
\]
Thus (GM26) makes the exponential independent of $c$ and proves
sufficiency.  For necessity, compare two tuples with the same last
$k-2$ roots, but first two sign pairs respectively $(++),(--)$
and $(+-),(-+)$.  Their sums are identical and the ratio of their
$Q_t$ values is $e^{8it\delta_\rho\gamma}$.  Constancy forces this
ratio to be one, which is exactly (GM26).  At these resonant times
the descended unit has at $\kappa_{a,b}$ the nonzero value
\[
 \exp\{tk(\delta_\rho^2-\gamma^2)
              +2it\delta_\rho\gamma(k-2a-2b)\}.
\]
Its unique interpolation polynomial of degree below $q$ specifies it
completely; it need not equal the different unit $E_t$.

If $m_0\geq2$, choose a tuple with two distinct roots, say the first
two.  Every element $p(s)$ has identical coefficients of $\tau_1$
and $\tau_2$, both $p'(\kappa_{a,b})$.  The respective coefficients
of $Q_t$ are its common nonzero constant coefficient multiplied by
$2t(\epsilon_1\delta_\rho+i\eta_1\gamma)$ and by
$2t(\epsilon_2\delta_\rho+i\eta_2\gamma)$.  They are unequal for
$t>0$, since these roots are distinct.  Those first coefficients
survive modulo $\tau_j^{m_0}$, proving nonmembership, even at the
resonant times of the simple-root case.  Finally $Q_tQ_{-t}=1$
in the finite algebra, and it commutes with $\mathsf U_h$.
This proves (GM27) with its stated inverse.
\end{proof}

For $N\geq q-1$, keep the original remainder map
$J_{\chi,N}:\mathbb C[S]_{\leq N}\to E_\chi$ in its power coordinates.
Its kernel is the complete ideal $\chi\mathbb C[S]_{\leq N-q}$.
Equation (GM25) shows that this kernel is unchanged under either
Gaussian heat: both physical maps are the original injective sum
embedding followed by units.  Hence the \emph{labelled} attained
forms are exactly
\begin{equation}\tag{GM28}
 H_\diamond(t)=
  (J_{\chi,N}\mathcal G_\diamond(t)^{-1}J_{\chi,N}^*)^{-1}
 \quad(\diamond={\rm fac},{\rm sum}),
\end{equation}
and GM-C--D compute their derivatives and remainders over every
coefficient of that ideal.  For sum heat, the metric at a fixed
physical vector $\mathsf U_h\iota(v)$ is
$E_t^{-*}H_{\rm sum}(t)E_t^{-1}$ in the original $v$ coordinates;
GM14--15 apply with $L_h$ replaced by multiplication by $(S-k/2)^2$
in the full algebra $E_\chi$.  For factor heat, the physical vector
with label $v$ is instead $\mathsf U_hQ_t\iota(v)$ and belongs to
the moving subspace in (GM27).  Outside the exact cases (GM26) a
congruence by a putative sum-coordinate unit is false.  The full
tensor map (GM27) is its correct replacement, and (GM28) remains the
attained metric of every such moving labelled vector.

\subsection{The precise relation to Connes's scalar Gaussian family}

For original coefficient vectors $c,d$ define the linear functional
on bounded functions of $y$ by
\[
 \mathfrak L_{c,d}(\varphi)
   =\int_{\mathbb R}\varphi(y)
          \overline{b_N(y)c}\,b_N(y)d\,w_h(y)\,dy.
\]
It is well-defined by Cauchy--Schwarz and (GM5), and the exact
Gaussian receiver is
\begin{equation}\tag{GM29}
 c^*G(t)d=\mathfrak L_{c,d}(\widehat F_{2t}),\qquad
 \widehat F_{2t}(y)=e^{-2ty^2}.
\end{equation}
The factor $2t$ is forced by the norm square.  Connes's zero-side
functional, instead, evaluates $\mathcal M f_t$ at zeta zeros:
its Gaussian summand at a zero $\rho$ is
$\mathcal M f_t(\rho)=e^{t(\rho-1/2)^2}$, with the original zero
multiplicity.  On a finite set of zeros this follows immediately
from (GM1); the full explicit-formula extension is Connes's Mellin
formula in Section 2 \cite{ConnesHeat}, with its archimedean and
prime terms retained.  His interpretation as a heat trace of a
self-adjoint operator with zeta ordinates assumes RH, as stated in
his Introduction and Theorem 1.1.  The original-source identities
(GM1)--(GM29) make no such assumption: they pair the same explicit
Gaussian transform with the explicitly different source functional.

The source heat $T(t)$ is not trace class for $t>0$.  In its unitary
Fourier realization choose disjoint measurable subsets of $[-1,1]$
of positive measure and their normalized indicators.  They form an
orthonormal sequence.  Their images under multiplication by
$e^{-ty^2}$ have pairwise disjoint supports and norms at least
$e^{-t}$.  Thus the images have no convergent subsequence, proving
that $T(t)$ is not compact; every trace-class operator is compact.
Consequently the source Gram in (GM29) is not a trace of this heat
operator and is not obtained by declaring a unitary identification
with Connes's discrete trace.  The exact transforms (GM2), (GM20)
and (GM29), and the finite jet maps (GM3), (GM25)--(GM28), specify
the actual relations used here.

For each specified finite $N,k$ and divisor, the constants in
(GM8), (GM15) and (GM22) are finite and computed from the full
original integrals and units.  For a finite determinant return,
(GM16) gives an attained error tolerance by choosing
$t\leq\varepsilon/(2\sum_\nu|a_\nu|r_\nu\gamma_{N_\nu})$
when that denominator is positive; a zero denominator gives zero
error.  This is a finite calculation and convergence statement.
No bound uniform in a growing degree, tensor order, packet size or
moving divisor is asserted.  In particular a fixed-order scalar
heat asymptotic from Connes is not substituted for growth estimates
on these actual moment matrices or their attained current receivers.
